\documentclass[a4paper,11pt]{article}
\usepackage[utf8]{inputenc}
\usepackage[T1]{fontenc}
\usepackage[french]{babel}
\usepackage[left=2.5cm,right=2.5cm,top=2cm,bottom=2cm]{geometry}
\usepackage{hyperref}
\usepackage{xcolor}
\usepackage{fontawesome5}
\usepackage{lmodern}

% --- Couleurs Dassault Systèmes ---
\definecolor{primary}{RGB}{0, 91, 154}

% --- Configuration des liens ---
\hypersetup{colorlinks=true, linkcolor=primary, urlcolor=primary}

\begin{document}

\pagestyle{empty}

% --- EN-TÊTE ---
\noindent
\begin{minipage}[t]{0.55\textwidth}
\textbf{Loïc Aron MBASSI EWOLO} \\
Élève Ingénieur 2A — Double Diplôme ENSEA \\
\faMapMarker\ Cergy, France \\
\faPhone\ (+33) 7 59 52 33 98 \\
\faEnvelope\ \href{mailto:loicaron.mbassiewolo@ensea.fr}{loicaron.mbassiewolo@ensea.fr} \\
\faLinkedin\ \href{https://www.linkedin.com/in/mbassi-loic-3942a9246/}{linkedin.com/in/mbassi-loic}
\end{minipage}
\hfill
\begin{minipage}[t]{0.40\textwidth}
\raggedleft
\textbf{Dassault Systèmes} \\
Campus de Vélizy-Villacoublay \\
Équipe 3DExperience Platform \\
Vélizy-Villacoublay, France
\end{minipage}

\vspace{1cm}

\noindent Cergy, le 23 décembre 2024

\vspace{0.8cm}

\noindent \textbf{Objet : Candidature au stage Ingénieur Développement IA (F/H)}

\vspace{0.6cm}

\noindent Madame, Monsieur,

\vspace{0.4cm}

Actuellement élève ingénieur en deuxième année à l'ENSEA dans le cadre d'un double diplôme avec l'École Polytechnique de Yaoundé, je recherche un stage de 4 à 6 mois à partir d'avril 2026. L'opportunité de contribuer au développement d'agents IA au sein de l'équipe 3DExperience correspond précisément à mon parcours et à mes ambitions professionnelles.

\vspace{0.3cm}

Le développement d'outils IA pour améliorer la qualité du développement logiciel est un domaine que j'ai déjà exploré concrètement. Mon expérience récente sur un \textbf{système multi-agents} (Google ADK, Gemini, LangChain) m'a permis de concevoir et orchestrer 5 agents IA spécialisés, en maîtrisant leur architecture, les interactions entre composants et le déploiement via Docker. Cette expérience m'a donné une compréhension approfondie du fonctionnement interne des agents IA, au-delà de leur simple utilisation.

\vspace{0.3cm}

En parallèle, j'ai développé \textbf{Laravel API Generator}, un outil open source de génération automatique de code à partir de diagrammes UML. Ce projet illustre ma sensibilité à la qualité du code et à l'automatisation des tâches répétitives — des préoccupations directement alignées avec votre mission d'amélioration de la qualité de développement.

\vspace{0.3cm}

Concernant les méthodologies TDD/BDD et les principes Clean Code, mon expérience chez \textbf{CSC} (plateforme Fintech) m'a formé aux pratiques rigoureuses : tests unitaires et d'intégration, CI/CD, revues de code systématiques et méthodologie Agile/Scrum. Je suis convaincu que la qualité logicielle passe par ces pratiques, et je serai enthousiaste à l'idée de les approfondir au sein de votre équipe.

\vspace{0.3cm}

Mon stage au \textbf{MILA} (Institut Québécois d'Intelligence Artificielle) en 2025 a renforcé ma maîtrise de Python et PyTorch dans un contexte de recherche exigeant, avec reproduction d'expériences SOTA et analyse rigoureuse des résultats. Cette expérience m'a appris à travailler de manière autonome tout en communiquant efficacement sur des sujets techniques complexes.

\vspace{0.3cm}

Intégrer Dassault Systèmes représente pour moi l'opportunité de contribuer à des solutions innovantes au sein d'une entreprise leader de l'innovation technologique. Je suis particulièrement motivé par la dimension collaborative de votre équipe et par l'impact concret des outils développés sur la plateforme 3DExperience.

\vspace{0.3cm}

Je reste disponible pour un entretien à votre convenance afin de vous présenter plus en détail mon parcours et ma motivation.

\vspace{0.4cm}

Dans l'attente de votre retour, je vous prie d'agréer, Madame, Monsieur, l'expression de mes salutations distinguées.

\vspace{1cm}

\noindent \textbf{Loïc Aron MBASSI EWOLO}

\end{document}
